\section{Introduction}

Most modern software architectures are distributed by design and message brokers play an important role in enabling communication and coordination between different parts of these systems. Using a message broker as an intermediary between components allows services to reliably send messages to each other without needing to know about the details of the recipient's implementation or availability. As a result, message brokers have become a critical component in many distributed systems and their performance can have a significant impact on the overall system's responsiveness and scalability. Among the many available message brokers, RabbitMQ is one of the most widely known and used open-source options. RabbitMQ's Quorum Queues, built on the Raft consensus algorithm~\cite{raft_paper}, are advertised as the go-to queue type for systems that require data safety and high availability but don't rely on time-sensitive messaging (such as stock tickers or instant messaging systems)~\cite{rmq_quorum_docs}.

However, the data safety advantages of Quorum Queues come at a cost. RabbitMQ can be deployed as a cluster of multiple nodes that together form a single logical broker. The Raft consensus algorithm requires every message to be replicated to a majority of these nodes before it is acknowledged, introducing a trade-off: scaling horizontally by adding nodes increases the total capacity and fault tolerance of the cluster, but it also increases the number of nodes that must agree on every write. The resulting overhead from additional network communication, disk writes and internal coordination makes it difficult to predict how throughput and latency will behave as the cluster grows.

As systems scale, understanding how RabbitMQ and its Quorum Queues perform under different load conditions becomes important. Therefore, this work aims to investigate this behavior by focusing on throughput and end-to-end latency as the primary performance metrics across varying cluster sizes and load conditions. By developing and using a custom benchmark CLI tool to benchmark RabbitMQ clusters, we evaluate how these two metrics change under horizontal scaling and identify any diminishing returns that may arise.

The remainder of this paper is structured as follows. Section~II provides background on RabbitMQ, Quorum Queues and the Raft consensus algorithm. Section~III describes the benchmark setup, including the infrastructure, the benchmark tool and the experiments. Section~IV analyzes the measurement results, whereas Section~V outlines directions for future work and Section~VI concludes the paper.
% DONE

